% SPDX-FileCopyrightText: 2026 David Shirvanyants
% SPDX-License-Identifier: CC0-1.0

\chapter{Introduction}

CLIP Slicer is a desktop application for loading, viewing, transforming, and
processing three-dimensional models for additive manufacturing. It presents
models in an interactive OpenGL view and keeps model transformations separate
from the source coordinates, allowing a model to be repositioned, rotated, or
scaled without modifying the loaded geometry.

The application currently supports two model formats:

\begin{description}
    \item[Binary STL] A triangulated surface model. CLIP Slicer reads triangle
        vertices and normals and displays the resulting mesh. Selected surface
        models can be combined and exported to binary STL after applying their
        document transformations.
    \item[CLI] A sliced model represented by layers and contours. CLIP Slicer
        reads binary CLI files and displays their layers as a three-dimensional
        model. It exports sliced models in CLI format.
\end{description}

CLIP Slicer uses a multiple-document interface. A document can contain one or
more STL and CLI models, and each model can be selected and shown or hidden
independently. Opening a file normally creates a new document. The
\emph{Open into document} command adds another model to the active document so
that several models can be viewed, positioned, and processed together.

Selected models are sliced on the build-platform layer grid defined by the
default layer thickness and absolute first-layer offset configured in the
application settings. The resulting sliced model can
be added to the current document or placed in a new document. When sliced
models are exported, the selected models can be combined in one CLI file. Their
layers are ordered by height, including layers from different models that
occupy the same height.

Unsupported-area analysis identifies surfaces which require additional print
support. Orientation optimization uses this analysis to rotate each selected
model toward an orientation with less unsupported area while retaining all
transformations separately from the source geometry. Support generation can
place contact tips and route tapered external pillars from those tips to bases
on the build platform.

The current export format is CLI. Support for additional three-dimensional and
layer formats, along with manual contour editing, is planned for later
development.
